\documentclass{article}
\usepackage{fullpage}
\begin{document}
\noindent\textbf{Describe your previous research experience or equivalent experience on complex projects, including the level of independence, while working as a member of a research/project team.}
\\

I am currently involved in three long-term research projects: two in computational biophysics, and one
in quantum algorithms. 
\\

As an undergraduate research assistant in the Vermaas lab at the MSU-DOE Plant Research Laboratory,
I am working on understanding the effects of alcohols on bacterial membrane structure, with the broader goal
of informing bioengineering efforts towards the design of tolerant strains for industrial biofuel production.
Over the last four months, I have conducted molecular dynamics (MD) simulations of bacterial
membranes with varying alcohol concentrations. I prepare simulation systems using Visual Molecular Dynamics (VMD), run MD simulations on high-performance computing clusters like MSU’s HPCC and NCSA
DeltaAI using Nanoscale Atomistic Molecular Dynamics (NAMD), and analyze the resulting trajectories by
computing membrane properties such as area per lipid and lipid order parameters
using Python and Tcl. I am currently writing a first-author manuscript for this work, using the results from
simulations to draw conclusions about membrane responses to solvent stress.
\\

I have also spent the last year working on characterizing water transport modulation mechanisms in
capped aquaporin channels. The aim of this project is to understand plant drought tolerance and responses,
and identify potential targets for increasing drought resilience. I set up and run MD simulations exposing the channel to different osmotic gradients, and analyze the resulting trajectories by computing pore radius profiles and water permeation rates. I am also responsible for producing figures for meetings with experimental collaborators, and conferences such as ACS Fall 2025 and the First Great Lakes Plant Science Conference.
\\

As an Honors Research Scholar, I am conducting research on quantum subgroup simulators under the
guidance of Professor Ryan LaRose. My work involves developing methods for efficiently simulating quantum
circuits on classical hardware, which is essential for the development of large-scale quantum
algorithms that cannot yet be run on existing quantum hardware. Over the last three months, I have explored
literature on group-theoretic and Lie-theoretic approaches to unitary subgroup membership, designed optimization routines to identify optimized gatesets for simulation, and created circuit rewriting
algorithms to reduce simulation costs. I discuss my progress weekly with Professor LaRose to reevaluate
future directions and ensure alignment with broader research goals.
\end{document}